\scanpage{38}\mspage{37}
% L1–4
de Boole en question est finie, alors les comp. connexes sont les \deleted{\unsure{$Y$|$X$}} $X\ne\varnothing$ \uline{minimaux} qui soient ouverts et fermés. [en supposant $0\ne1$
% L5–7
i.e. $M$ non réduit à un p\textsuperscript{t} --- sinon on a \deleted{une} \unsure{mor\textsuperscript{t}|moral\textsuperscript{t}|ess\textsuperscript{t}} une structure ``vide'' \ldots\ ]

% L8–10
\uline{Ex}\qquad Reprenons $M=\mathrm{Mod}(\unsure{L|\mathcal L})$, et $M_A$ où $A\in M=\mathrm{Mod}(L)$, donc $A$ partie modérée d'un tronçon.

% L11–12
On aura donc
\[
\begin{aligned}
A=X_{\unsure{\varphi|\Phi}} &,\quad
\varphi\in{\unsure{\mathrm{Drop}|\mathrm{Drp}}}^{\unsure{*|K}}
  \bigl(\unsure{\mathrm{Tro}|\mathrm{Tr}_0}(L)\bigr)\\[-0.6ex]
 &\mkern-8mu\inserted{=\bigcup_i T_i}
\end{aligned}
\]
% L13–14
Je dis que les comp. connexes de $A$ sont les \unsure{ps.tronçons|pr.tronçons} composant $T_i$.
% L15–17
En effet, on a bien \unsure{sur|sûr}
\[
A=\bigvee_i T_i\quad\text{dans }M,\quad\deleted{\unsure{\&}}\quad\text{et}\quad T_i\wedge T_j=\varnothing
\]
% L18–22
si $i\ne j$, il suffit donc de prouver que a) Chaque $T_i$ est ouvert et fermé dans $A$, et b) Chaque $T_i$ est connexe, i.e. ne peut se décomposer en deux parties qui \deleted{\unsure{soit|soient}} disjointes et \inserted{rel\textsuperscript{t}} fermées $\ne\varnothing$.

